% !PRINTED 190
\weberpage{190}

Wenn wir nämlich unter $\xi_1, \xi_2$ neue Variable verstehen und
\weq{8}{x_1' = \sqrt{2}\,\xi_1\,\xi_2,\quad x_2' = \xi_1^2,\quad x_3' = \xi_2^2}
setzen, so verschwindet $x_1'^2 - 2\,x_2'\,x_3'$ identisch und
$y_1', y_2', y_3'$ gehen durch (5) und (8) in binäre quadratische
Formen der Variablen $\xi_1, \xi_2$ über, die nach (7) der
Gleichung
\weq{9}{y_1'^2 - 2\,y_2'\,y_3' = 0}
identisch genügen müssen, d. h. $y_2' y_3'$ muss ein
vollständiges Quadrat werden. Die quadratischen Formen $y_1', y_2',
y_3'$ zerlegen wir nun in je zwei lineare Factoren. Dabei können
$y_2'$ und $y_3'$ keinen gemeinsamen Factor haben, da sonst auch
der andere Factor in beiden Functionen übereinstimmen und also
die drei Co\"efficienten von $y_2'$ und $y_3'$ mit einander
proportional sein müssten. Dann würde aber $|A'|$ verschwinden,
was nicht möglich ist. Demnach ergiebt sich aus (9), dass $y_2'$
und $y_3'$ Quadrate linearer Functionen sein müssen. Setzen wir
\weq{10}{\eta_1 = \alpha\,\xi_1 + \beta\,\xi_2,\quad \eta_2 = \gamma\,\xi_1 + \delta\,\xi_2,}
so können wir also hiernach mit Rücksicht auf (9)
\weq{11}{y_1' = \sqrt{2}\,\eta_1\,\eta_2,\quad y_2' = \eta_1^2,\quad y_3' = \eta_2^2}
setzen; wenn wir die Multiplicationen ausführen und dann die
Substitution (8) im umgekehrten Sinne ausführen, so erhalten wir
aus (10) und (11):
\[
\begin{aligned}
y_1' &= (\alpha\delta + \beta\gamma)\,x_1' + \sqrt{2}\,\alpha\gamma\,x_2' + \sqrt{2}\,\beta\delta\,x_3', \\
y_2' &= \sqrt{2}\,\alpha\beta\,x_1' + \alpha^2\,x_2' + \beta^2\,x_3', \\
y_3' &= \sqrt{2}\,\gamma\delta\,x_1' + \gamma^2\,x_2' + \delta^2\,x_3'.
\end{aligned}
\]
Also lautet die Substitution $A'$:
\weq{12}{A' = \begin{pmatrix}
\alpha\delta + \beta\gamma, & \sqrt{2}\,\alpha\gamma, & \sqrt{2}\,\beta\delta \\
\sqrt{2}\,\alpha\beta, & \alpha^2, & \beta^2 \\
\sqrt{2}\,\gamma\delta, & \gamma^2, & \delta^2
\end{pmatrix}.}
Durch diese Substitution ist die Bedingung (7) noch nicht völlig
befriedigt, sondern es folgt nur, dass die rechte und linke Seite
sich durch einen constanten Factor unterscheiden, der von den
Co\"efficienten $\alpha, \beta, \gamma, \delta$ abhängt. Nun ist
aber der Co\"efficient von $x_1'^2$ in der Verbindung
$y_1'^2 - 2\,y_2'\,y_3'$:
\[
(\alpha\delta + \beta\gamma)^2 - 4\,\alpha\delta\beta\gamma = (\alpha\delta - \beta\gamma)^2,
\]
%CONT>>
